\noindent\begin{minipage}[t][\textheight][s]{\textwidth}

% --- Image pleine largeur ---
\illusloflabel{Kalam}%
\noindent\bwimage[width=\linewidth]{acte_05/scene_9/_0330_kalams-poudre-de-riz.jpg}\\[0.4em]
{\centering\figcaptionname\itshape Représentation rituelle de Bhadrakali à trente-deux bras, tracée au sol avec des poudres naturelles dans le cadre du Vadakkupurathu Pattu\par}

\vspace{0.2\baselineskip}

% --- Texte d'explication encadré ---
\noindent\fbox{%
  \begin{minipage}{\dimexpr\linewidth-2\fboxsep-2\fboxrule\relax}
  \vspace{0.3\baselineskip}
  \fontsize{9}{11}\selectfont
  \begin{description}
    \setlength{\itemsep}{0.5\baselineskip}
    \item[\textbf{Bhadrakali}] — déesse hindoue de la tradition kéralaise, forme féroce de Kali/Parvati, victorieuse du démon Darika. Divinité centrale du Kerala.
    \item[\textbf{kalam}] — dessin rituel tracé au sol avec des poudres de couleurs naturelles (farine de riz, charbon, curcuma, chaux, feuilles broyées). \textit{Kalam} signifie «~image~» ou «~dessin~», \textit{ezhuthu} «~écrire/tracer~».
    \item[\textbf{à trente-deux bras}] — le nombre de bras n'est pas anecdotique : c'est le nombre de bras de la déesse qui détermine les dimensions du kalam. Dans le Vadakkupurathu Pattu, le rituel progresse sur 12~jours : la déesse est représentée avec 8~bras (jours~1--4), puis 16~bras (jours~5--8), puis 32~bras (jours~9--11), et enfin 64~bras le dernier jour, dans sa forme la plus puissante, montée sur Vetalaroopa.
    \item[\textbf{Vadakkupurathu Pattu}] — littéralement «~le chant/rituel du côté nord~» (\textit{vadakku}~= nord, \textit{puram}~= côté extérieur, \textit{pattu}~= chant). C'est un \textit{Kalamezhuthu Pattu} de 12~jours célébré une fois tous les 12~ans pour apaiser la déesse Bhadrakali, dans la cour nord du temple Vaikom Mahadeva (Kerala). Il est accompagné de rituels spéciaux et d'offrandes \textit{Guruthi}.
  \end{description}
  \vspace{0.3\baselineskip}
  \end{minipage}%
}

\vspace*{\fill}

\end{minipage}
